% Options for packages loaded elsewhere
\PassOptionsToPackage{unicode}{hyperref}
\PassOptionsToPackage{hyphens}{url}
\documentclass[
  ignorenonframetext,
]{beamer}
\newif\ifbibliography
\usepackage{pgfpages}
\setbeamertemplate{caption}[numbered]
\setbeamertemplate{caption label separator}{: }
\setbeamercolor{caption name}{fg=normal text.fg}
\beamertemplatenavigationsymbolsempty
% remove section numbering
\setbeamertemplate{part page}{
  \centering
  \begin{beamercolorbox}[sep=16pt,center]{part title}
    \usebeamerfont{part title}\insertpart\par
  \end{beamercolorbox}
}
\setbeamertemplate{section page}{
  \centering
  \begin{beamercolorbox}[sep=12pt,center]{section title}
    \usebeamerfont{section title}\insertsection\par
  \end{beamercolorbox}
}
\setbeamertemplate{subsection page}{
  \centering
  \begin{beamercolorbox}[sep=8pt,center]{subsection title}
    \usebeamerfont{subsection title}\insertsubsection\par
  \end{beamercolorbox}
}
% Prevent slide breaks in the middle of a paragraph
\widowpenalties 1 10000
\raggedbottom
\AtBeginPart{
  \frame{\partpage}
}
\AtBeginSection{
  \ifbibliography
  \else
    \frame{\sectionpage}
  \fi
}
\AtBeginSubsection{
  \frame{\subsectionpage}
}
\usepackage{iftex}
\ifPDFTeX
  \usepackage[T1]{fontenc}
  \usepackage[utf8]{inputenc}
  \usepackage{textcomp} % provide euro and other symbols
\else % if luatex or xetex
  \usepackage{unicode-math} % this also loads fontspec
  \defaultfontfeatures{Scale=MatchLowercase}
  \defaultfontfeatures[\rmfamily]{Ligatures=TeX,Scale=1}
\fi
\usepackage{lmodern}
\ifPDFTeX\else
  % xetex/luatex font selection
\fi
% Use upquote if available, for straight quotes in verbatim environments
\IfFileExists{upquote.sty}{\usepackage{upquote}}{}
\IfFileExists{microtype.sty}{% use microtype if available
  \usepackage[]{microtype}
  \UseMicrotypeSet[protrusion]{basicmath} % disable protrusion for tt fonts
}{}
\makeatletter
\@ifundefined{KOMAClassName}{% if non-KOMA class
  \IfFileExists{parskip.sty}{%
    \usepackage{parskip}
  }{% else
    \setlength{\parindent}{0pt}
    \setlength{\parskip}{6pt plus 2pt minus 1pt}}
}{% if KOMA class
  \KOMAoptions{parskip=half}}
\makeatother
\usepackage{color}
\usepackage{fancyvrb}
\newcommand{\VerbBar}{|}
\newcommand{\VERB}{\Verb[commandchars=\\\{\}]}
\DefineVerbatimEnvironment{Highlighting}{Verbatim}{commandchars=\\\{\}}
% Add ',fontsize=\small' for more characters per line
\newenvironment{Shaded}{}{}
\newcommand{\AlertTok}[1]{\textcolor[rgb]{1.00,0.00,0.00}{\textbf{#1}}}
\newcommand{\AnnotationTok}[1]{\textcolor[rgb]{0.38,0.63,0.69}{\textbf{\textit{#1}}}}
\newcommand{\AttributeTok}[1]{\textcolor[rgb]{0.49,0.56,0.16}{#1}}
\newcommand{\BaseNTok}[1]{\textcolor[rgb]{0.25,0.63,0.44}{#1}}
\newcommand{\BuiltInTok}[1]{\textcolor[rgb]{0.00,0.50,0.00}{#1}}
\newcommand{\CharTok}[1]{\textcolor[rgb]{0.25,0.44,0.63}{#1}}
\newcommand{\CommentTok}[1]{\textcolor[rgb]{0.38,0.63,0.69}{\textit{#1}}}
\newcommand{\CommentVarTok}[1]{\textcolor[rgb]{0.38,0.63,0.69}{\textbf{\textit{#1}}}}
\newcommand{\ConstantTok}[1]{\textcolor[rgb]{0.53,0.00,0.00}{#1}}
\newcommand{\ControlFlowTok}[1]{\textcolor[rgb]{0.00,0.44,0.13}{\textbf{#1}}}
\newcommand{\DataTypeTok}[1]{\textcolor[rgb]{0.56,0.13,0.00}{#1}}
\newcommand{\DecValTok}[1]{\textcolor[rgb]{0.25,0.63,0.44}{#1}}
\newcommand{\DocumentationTok}[1]{\textcolor[rgb]{0.73,0.13,0.13}{\textit{#1}}}
\newcommand{\ErrorTok}[1]{\textcolor[rgb]{1.00,0.00,0.00}{\textbf{#1}}}
\newcommand{\ExtensionTok}[1]{#1}
\newcommand{\FloatTok}[1]{\textcolor[rgb]{0.25,0.63,0.44}{#1}}
\newcommand{\FunctionTok}[1]{\textcolor[rgb]{0.02,0.16,0.49}{#1}}
\newcommand{\ImportTok}[1]{\textcolor[rgb]{0.00,0.50,0.00}{\textbf{#1}}}
\newcommand{\InformationTok}[1]{\textcolor[rgb]{0.38,0.63,0.69}{\textbf{\textit{#1}}}}
\newcommand{\KeywordTok}[1]{\textcolor[rgb]{0.00,0.44,0.13}{\textbf{#1}}}
\newcommand{\NormalTok}[1]{#1}
\newcommand{\OperatorTok}[1]{\textcolor[rgb]{0.40,0.40,0.40}{#1}}
\newcommand{\OtherTok}[1]{\textcolor[rgb]{0.00,0.44,0.13}{#1}}
\newcommand{\PreprocessorTok}[1]{\textcolor[rgb]{0.74,0.48,0.00}{#1}}
\newcommand{\RegionMarkerTok}[1]{#1}
\newcommand{\SpecialCharTok}[1]{\textcolor[rgb]{0.25,0.44,0.63}{#1}}
\newcommand{\SpecialStringTok}[1]{\textcolor[rgb]{0.73,0.40,0.53}{#1}}
\newcommand{\StringTok}[1]{\textcolor[rgb]{0.25,0.44,0.63}{#1}}
\newcommand{\VariableTok}[1]{\textcolor[rgb]{0.10,0.09,0.49}{#1}}
\newcommand{\VerbatimStringTok}[1]{\textcolor[rgb]{0.25,0.44,0.63}{#1}}
\newcommand{\WarningTok}[1]{\textcolor[rgb]{0.38,0.63,0.69}{\textbf{\textit{#1}}}}
\setlength{\emergencystretch}{3em} % prevent overfull lines
\providecommand{\tightlist}{%
  \setlength{\itemsep}{0pt}\setlength{\parskip}{0pt}}
% Slide layout defaults for warning-clean scientific decks.
\usepackage{etoolbox}
\IfFileExists{xurl.sty}{\usepackage{xurl}}{}
\IfFileExists{seqsplit.sty}{\usepackage{seqsplit}}{\newcommand{\seqsplit}[1]{#1}}
\protected\def\breakseq#1{\seqsplit{#1}}
\protected\def\breaktt#1{\begingroup\ttfamily\seqsplit{#1}\endgroup}
\setlength{\emergencystretch}{6em}
\tolerance=5000
\hbadness=10000
\hfuzz=1pt
\setlength{\tabcolsep}{2pt}
\AtBeginEnvironment{longtable}{\tiny\renewcommand{\arraystretch}{0.86}\setlength{\tabcolsep}{1pt}}
\AtBeginEnvironment{tabular}{\tiny\renewcommand{\arraystretch}{0.86}\setlength{\tabcolsep}{1pt}}

% Natbib and cross-reference fallbacks — slides don't load natbib
% or cleveref, but combined-PDF manuscript prose may emit these
% commands. The fallback renders citations as a bracketed key list
% and cross-references as detokenized labels so slides don't fail on
% undefined control sequences or raw underscores. \providecommand is
% a no-op if packages load later.
\providecommand{\citep}[1]{[#1]}
\providecommand{\citet}[1]{#1}
\providecommand{\citealp}[1]{#1}
\providecommand{\citeauthor}[1]{#1}
\providecommand{\citeyear}[1]{#1}
\providecommand{\cref}[1]{\texttt{\detokenize{#1}}}
\providecommand{\Cref}[1]{\texttt{\detokenize{#1}}}
\usepackage{bookmark}
\IfFileExists{xurl.sty}{\usepackage{xurl}}{} % add URL line breaks if available
\urlstyle{same}
% fallback for those not using the hyperref driver hyperxmp:
\makeatletter
\@ifundefined{xmpquote}{\newcommand{\xmpquote}[1]{#1}}{}
\makeatother
\hypersetup{
  hidelinks,
  pdfcreator={LaTeX via pandoc}}

\author{\texorpdfstring{}{}}
\date{}

\begin{document}

\section{Pipeline Internals}\label{sec:pipeline_internals}

\begin{frame}[allowframebreaks]{Pipeline Internals}
This supplemental section documents the data structures and on-disk
artifacts the pipeline produces, for readers who want to extend or audit
it.
\end{frame}

\begin{frame}[fragile,allowframebreaks]{Data structures}
\protect\phantomsection\label{data-structures}
The Mermaid class diagram in this subsection shows the canonical fields
each record carries through the pipeline. Records have additional
optional metadata (e.g.~\texttt{Paper.url}, \texttt{Paper.publisher},
\texttt{Paper.isbn}, \texttt{Paper.raw}) omitted for readability ---
consult \breaktt{infrastructure/search/literature/models.py}
(\href{https://github.com/docxology/template/tree/main/infrastructure/search/literature}{source
on GitHub}) and \texttt{src/pipeline.py}
(\href{https://github.com/docxology/template/tree/main/projects/templates/template_search_project/src}{source
on GitHub}) for the full schema.

\begin{Shaded}
\begin{Highlighting}[]
\NormalTok{classDiagram}
\NormalTok{    class SearchQuery \{}
\NormalTok{        +str text}
\NormalTok{        +int max\_results}
\NormalTok{        +int|None year\_min}
\NormalTok{        +int|None year\_max}
\NormalTok{        +list\textasciitilde{}str\textasciitilde{} sources}
\NormalTok{        +list\textasciitilde{}str\textasciitilde{} fields}
\NormalTok{        +matches\_year(year) bool}
\NormalTok{    \}}

\NormalTok{    class Paper \{}
\NormalTok{        +str id}
\NormalTok{        +str title}
\NormalTok{        +list\textasciitilde{}str\textasciitilde{} authors}
\NormalTok{        +int|None year}
\NormalTok{        +str|None doi}
\NormalTok{        +str|None url}
\NormalTok{        +str|None abstract}
\NormalTok{        +str|None pdf\_url}
\NormalTok{        +str|None fulltext}
\NormalTok{        +str|None venue}
\NormalTok{        +str|None venue\_type}
\NormalTok{        +str|None volume}
\NormalTok{        +str|None issue}
\NormalTok{        +str|None pages}
\NormalTok{        +str|None publisher}
\NormalTok{        +str|None edition}
\NormalTok{        +str|None isbn}
\NormalTok{        +list\textasciitilde{}str\textasciitilde{} keywords}
\NormalTok{        +str|None source}
\NormalTok{        +float score}
\NormalTok{        +dict raw}
\NormalTok{        +to\_dict() dict}
\NormalTok{        +from\_dict(d) Paper}
\NormalTok{    \}}

\NormalTok{    class SearchResult \{}
\NormalTok{        +SearchQuery query}
\NormalTok{        +list\textasciitilde{}Paper\textasciitilde{} papers}
\NormalTok{        +dict per\_source\_counts}
\NormalTok{        +dict errors}
\NormalTok{        +to\_dict() dict}
\NormalTok{        +to\_json() str}
\NormalTok{    \}}

\NormalTok{    class BibEntry \{}
\NormalTok{        +str entry\_type}
\NormalTok{        +str citation\_key}
\NormalTok{        +OrderedDict fields}
\NormalTok{    \}}

\NormalTok{    class BibDatabase \{}
\NormalTok{        +list\textasciitilde{}BibEntry\textasciitilde{} entries}
\NormalTok{        +str|None preamble}
\NormalTok{        +find(key) BibEntry}
\NormalTok{    \}}

\NormalTok{    class LiteratureRunArtifacts \{}
\NormalTok{        +SearchResult result}
\NormalTok{        +Path|None corpus\_path}
\NormalTok{        +Path|None bibtex\_path}
\NormalTok{        +list\textasciitilde{}FetchResult\textasciitilde{} enrichment\_log}
\NormalTok{        +Path|None cache\_dir}
\NormalTok{        +dict citation\_keys}
\NormalTok{        +papers() list\textasciitilde{}Paper\textasciitilde{}}
\NormalTok{    \}}

\NormalTok{    class ManuscriptVariables \{}
\NormalTok{        +str config\_query}
\NormalTok{        +int config\_max\_results}
\NormalTok{        +str config\_sources}
\NormalTok{        +int result\_num\_papers}
\NormalTok{        +int result\_num\_sources}
\NormalTok{        +str result\_per\_source}
\NormalTok{        +str result\_errors}
\NormalTok{        +str result\_year\_min}
\NormalTok{        +str result\_year\_max}
\NormalTok{        +int result\_with\_abstract}
\NormalTok{        +int result\_with\_doi}
\NormalTok{        +int deep\_max\_results\_per\_keyword}
\NormalTok{        +int deep\_keyword\_count}
\NormalTok{        +str deep\_keywords\_joined}
\NormalTok{        +str deep\_sources}
\NormalTok{        +str deep\_unique\_papers}
\NormalTok{        +as\_dict() dict}
\NormalTok{        +as\_uppercase\_keys() dict}
\NormalTok{    \}}

\NormalTok{    class SynthesisResult \{}
\NormalTok{        +str kind}
\NormalTok{        +str prompt}
\NormalTok{        +str text}
\NormalTok{        +str|None paper\_id}
\NormalTok{    \}}

\NormalTok{    SearchResult {-}{-}\textgreater{} SearchQuery}
\NormalTok{    SearchResult {-}{-}\textgreater{} "*" Paper}
\NormalTok{    LiteratureRunArtifacts {-}{-}\textgreater{} SearchResult}
\NormalTok{    LiteratureRunArtifacts {-}{-}\textgreater{} "*" BibEntry : via paper\_to\_bibentry}
\NormalTok{    BibDatabase {-}{-}\textgreater{} "*" BibEntry}
\end{Highlighting}
\end{Shaded}
\end{frame}

\begin{frame}[fragile,allowframebreaks]{On-disk layout}
\protect\phantomsection\label{on-disk-layout}
The project keeps committed input data in \texttt{data/}, regeneratable
outputs in \texttt{output/} (gitignored), and the manuscript source in
\texttt{manuscript/}. The Mermaid flowchart in this subsection (rendered
in the HTML build; the PDF build strips Mermaid) lists every artefact
the standard pipeline writes:

\begin{itemize}
\tightlist
\item
  \breaktt{data/corpus.json} --- the bundled offline corpus, CI-safe
  default for \texttt{LocalBackend}.
\item
  \breaktt{output/search/results.json} --- raw \texttt{SearchResult} JSON
  from the latest run.
\item
  \texttt{output/search/cache/search\_\textless{}hash\textgreater{}.json}
  --- deterministic \texttt{SearchCache} files.
\item
  \texttt{output/cache/abs/\textless{}safe\_id\textgreater{}.txt} ---
  one cached abstract per paper.
\item
  \texttt{output/cache/pdf/\textless{}safe\_id\textgreater{}.\{pdf,txt\}}
  --- PDF bytes plus extracted text.
\item
  \breaktt{output/llm/synthesis.md} --- corpus-level LLM synthesis (when
  enabled).
\item
  \texttt{output/llm/per\_paper/\textless{}safe\_id\textgreater{}.md}
  --- one per-paper note per paper (when enabled).
\item
  \texttt{output/figures/\{papers\_per\_source,year\_histogram,score\_distribution\}.png}
  --- diagnostic figures.
\item
  \breaktt{output/data/manuscript\_variables.json} --- substitution table
  consumed by the resolver.
\item
  \breaktt{output/corpus.json} --- enriched corpus, written in
  \texttt{LocalBackend}-compatible format so it can re-seed a
  deterministic future run.
\item
  \breaktt{output/enrichment\_log.json} --- one entry per fetcher per
  paper.
\item
  \breaktt{output/reading\_report.md} --- final markdown reading report.
\item
  \breaktt{output/run\_summary.json} --- one-line metadata for the run.
\item
  \breaktt{manuscript/references.bib} --- auto-populated, Pandoc-ready
  BibTeX.
\end{itemize}

\begin{Shaded}
\begin{Highlighting}[]
\NormalTok{flowchart TB}
\NormalTok{    P["projects/templates/template\_search\_project/"]}
\NormalTok{    P {-}{-}\textgreater{} DAT["data/\textless{}br/\textgreater{}committed"]}
\NormalTok{    P {-}{-}\textgreater{} OUT["output/\textless{}br/\textgreater{}regeneratable · gitignored"]}
\NormalTok{    P {-}{-}\textgreater{} MAN["manuscript/"]}

\NormalTok{    DAT {-}{-}\textgreater{} CORP\_IN[corpus.json\textless{}br/\textgreater{}bundled offline corpus\textless{}br/\textgreater{}CI{-}safe default]}

\NormalTok{    OUT {-}{-}\textgreater{} SR["search/"]}
\NormalTok{    OUT {-}{-}\textgreater{} CC["cache/"]}
\NormalTok{    OUT {-}{-}\textgreater{} LD["llm/"]}
\NormalTok{    OUT {-}{-}\textgreater{} FIG["figures/"]}
\NormalTok{    OUT {-}{-}\textgreater{} DD["data/"]}
\NormalTok{    OUT {-}{-}\textgreater{} CORP\_OUT[corpus.json\textless{}br/\textgreater{}enriched · LocalBackend{-}compatible]}
\NormalTok{    OUT {-}{-}\textgreater{} ELOG[enrichment\_log.json\textless{}br/\textgreater{}per{-}fetch status]}
\NormalTok{    OUT {-}{-}\textgreater{} RR[reading\_report.md]}
\NormalTok{    OUT {-}{-}\textgreater{} RS[run\_summary.json]}

\NormalTok{    SR {-}{-}\textgreater{} SR\_R[results.json\textless{}br/\textgreater{}SearchResult]}
\NormalTok{    SR {-}{-}\textgreater{} SR\_C["cache/search\_HASH.json\textless{}br/\textgreater{}deterministic SearchCache"]}

\NormalTok{    CC {-}{-}\textgreater{} CC\_A["abs/SAFE\_ID.txt"]}
\NormalTok{    CC {-}{-}\textgreater{} CC\_P["pdf/SAFE\_ID.pdf and .txt"]}

\NormalTok{    LD {-}{-}\textgreater{} LD\_S[synthesis.md]}
\NormalTok{    LD {-}{-}\textgreater{} LD\_PP["per\_paper/SAFE\_ID.md"]}

\NormalTok{    FIG {-}{-}\textgreater{} F1[papers\_per\_source.png]}
\NormalTok{    FIG {-}{-}\textgreater{} F2[year\_histogram.png]}
\NormalTok{    FIG {-}{-}\textgreater{} F3[score\_distribution.png]}

\NormalTok{    DD {-}{-}\textgreater{} MV[manuscript\_variables.json]}

\NormalTok{    MAN {-}{-}\textgreater{} BIB[references.bib\textless{}br/\textgreater{}auto{-}populated · Pandoc{-}ready]}

\NormalTok{    classDef dir fill:\#0f172a,stroke:\#0f172a,color:\#fff}
\NormalTok{    classDef gen fill:\#0f766e,stroke:\#0f172a,color:\#fff}
\NormalTok{    classDef src fill:\#1e3a8a,stroke:\#0f172a,color:\#fff}
\NormalTok{    class P,DAT,OUT,MAN,SR,CC,LD,FIG,DD dir}
\NormalTok{    class CORP\_OUT,ELOG,RR,RS,SR\_R,SR\_C,CC\_A,CC\_P,LD\_S,LD\_PP,F1,F2,F3,MV,BIB gen}
\NormalTok{    class CORP\_IN src}
\end{Highlighting}
\end{Shaded}
\end{frame}

\begin{frame}[fragile,allowframebreaks]{Citation-key collision handling}
\protect\phantomsection\label{citation-key-collision-handling}
\breaktt{paper\_to\_bibentry()} generates citation keys as
\texttt{\textless{}author\textgreater{}\textless{}year\textgreater{}\textless{}title-word\textgreater{}}
(with stop-words filtered and unicode folded). When two papers in the
same result set produce the same key --- common when one author
publishes multiple papers in the same year on closely related topics ---
\breaktt{src/pipeline.py::\_disambiguate\_citation\_key} appends a
deterministic suffix from the alphabet (\texttt{a}, \texttt{b}, \ldots,
\texttt{z}, then two-letter combinations \texttt{aa}, \texttt{ab},
\ldots) until uniqueness is restored, with a numeric \texttt{\_1},
\texttt{\_2}, \ldots{} fallback for the pathological case. The mapping
is exposed to downstream stages via
\breaktt{LiteratureRunArtifacts.citation\_keys}, and the report uses
these keys verbatim, so the LLM synthesis and the BibTeX file always
agree.

The deep-search workflow has its own collision handler in
\breaktt{src/deep\_search.py::run\_deep\_search} that operates over the
post-deduplication aggregate roster --- see \breakseq{[@sec:deep\_search]} ---
and the unified \breaktt{references\_deep.bib} reflects the same mapping.
\end{frame}

\begin{frame}[fragile,allowframebreaks]{Failure isolation}
\protect\phantomsection\label{failure-isolation}
\begin{itemize}
\tightlist
\item
  \textbf{A backend is unreachable.} \breaktt{LiteratureClient} records
  the message into \texttt{result.errors{[}name{]}} and continues. The
  reading report surfaces these errors in a callout block at the top.
\item
  \textbf{An abstract fetch fails.} The fetcher records
  \texttt{status="error"} in \breaktt{enrichment\_log.json} and the paper
  keeps its existing (possibly empty) abstract.
\item
  \textbf{A PDF fetch fails or \texttt{pypdf} is missing.} Same pattern
  --- the PDF, if downloaded, is still cached on disk. The reading
  report does not reference the missing fulltext.
\item
  \textbf{The LLM is unreachable or unconfigured.}
  \breaktt{scripts/run\_search\_pipeline.py} logs a warning and skips the
  synthesis stage entirely --- \texttt{output/llm/} is left empty and
  the reading report omits both the per-paper-notes and cross-corpus
  sections. No placeholder text is ever written into the archive, so a
  missing LLM is observable from the absence of those sections rather
  than from a fake ``(LLM unavailable)'' string.
\end{itemize}
\end{frame}

\end{document}
